% Complete quantitative analytic core; included by the main manuscript.
\section{Finite analytic constructions used below}
The complex notation in this section abbreviates pairs of rational
approximation procedures with certified errors. It does not quantify over
arbitrary real numbers or arbitrary functions. All final inputs $x,y$ are
positive rationals, all vertical contour sides used in the final estimate
are rational, and every infinite expression below is accompanied by a
tail bound. The proof only needs finite sums, products, powers, rational
division with a supplied positive denominator bound, and rational
approximations to the elementary exponential, sine, cosine and logarithm.
The latter can be obtained from their power series after range reduction.
For example, $\log u$ on a positive rational interval follows from
$\log u=2\sum_{j\ge0}z^{2j+1}/(2j+1)$, $z=(u-1)/(u+1)$;
binary range reduction puts $|z|\le1/3$, and the geometric tail is
explicit. Thus every expression used here has a primitive recursive
approximation at a requested precision.

\paragraph{Finite Euler summation.}
Let $s=\sigma+it$, $\sigma>0$, and let $1\le N<M$ be integers.
Integration by parts on each interval $[n,n+1]$ and telescoping gives
\[
 \sum_{n\le M}n^{-s}=\sum_{n\le N}n^{-s}
 +\frac{N^{1-s}-M^{1-s}}{s-1}
 -s\int_N^M\{u\}u^{-s-1}\,du.                              \tag{E1}
\]
For $s=1$ the quotient means its explicit logarithmic limit; later we
separate this point algebraically. One can verify (E1) without integration
notation: on $[n,n+1]$, replace $\{u\}$ by $u-n$, use the elementary
antiderivatives of $u^{-s}$ and $u^{-s-1}$, and sum. Both endpoint sums
telescope. The displayed integrals are therefore finite elementary
expressions before any limiting operation.

For $\sigma>1$ the term $M^{1-s}$ and the tail integral have explicit
limits, so (E1) agrees with the Dirichlet series definition of $\zeta$.
For $\sigma>0$, $s\ne1$, define instead
\[
 Z(s)=\sum_{n\le N}n^{-s}+\frac{N^{1-s}}{s-1}
       -s\int_N^\infty\{u\}u^{-s-1}\,du.                   \tag{E2}
\]
Equation (E1) proves that this does not depend on $N$. The tail after $V$
is at most $|s|V^{-\sigma}/\sigma$. This is the only continuation of
$\zeta$ used in this proof; no analytic-continuation uniqueness theorem is
needed. Differentiating the finite expression and its explicitly bounded
tails is legitimate, since for $j\ge0$
\[
 \int_V^\infty (\log u)^j u^{-\sigma-1}\,du
 =V^{-\sigma}\sum_{r=0}^j
     \frac{j!}{(j-r)!}\frac{(\log V)^{j-r}}{\sigma^{r+1}}.    \tag{E3}
\]
This identity follows by $j$ integrations by parts; in the application
$j\le3$ and $\sigma\ge1/2$. Taking $V=2^v$ with a sufficiently large
explicit integer $v$ makes each tail as small as requested. For instance,
one may start with $v=2^{q+10}(T+4)^4$ when $|t|\le T$ and the desired
error is $2^{-q}$; the exponential decay $2^{-v/2}$ dominates the degree
three polynomial in $v$ appearing in (E3). Every such comparison reduces
to the elementary inequalities $n\le2^n$ and finite multiplication.

Writing $\zeta$ for $Z$, the $N=1$ formula is
\[
 \zeta(s)=\frac1{s-1}+h(s),\qquad
 h(s)=1-s\int_1^\infty\{u\}u^{-s-1}\,du.                   \tag{E4}
\]
Formula (E3) proves directly that $h$ has three controlled complex
derivatives on every rectangle used below. The pole at $1$ is therefore
isolated by an explicit algebraic expression, rather than detected by a
zero or pole enumeration.

\paragraph{Finite Euler products and their tails.}
For $\sigma=1+\eta>1$ put
$P_N(s)=\prod_{p\le N}(1-p^{-s})^{-1}$. Every factor has a supplied
apartness bound $|1-p^{-s}|\ge1-p^{-\sigma}\ge1/2$.
Expand each factor as a geometric series. Unique prime factorization
identifies the resulting product with the sum of $n^{-s}$ over integers
whose prime factors are at most $N$. This assertion can be checked first
with all exponents bounded and then with geometric tails; no infinite
choice is involved. Every integer $n\le N$ has all its prime factors at
most $N$, and consequently
\[
        |\zeta(s)-P_N(s)|\le \sum_{n>N}n^{-\sigma}
                    \le N^{-\eta}/\eta.                   \tag{E5}
\]
The last inequality is comparison with $\int_N^\infty u^{-1-\eta}du$.
If $\eta\ge1/B$, the explicit choice
$N=2^{B(q+B+2)}$ makes the right-hand side at most $2^{-q-2}$:
$N^{-\eta}/\eta\le B2^{-q-B-2}\le2^{-q-2}$.
For a finite verification of the geometric expansion, truncate each prime
exponent at $J\ge q+2N+10$. The product error is bounded by
$N2^N2^{-J}$, because there are at most $N$ factors, each complete factor
has modulus at most $2$, and each geometric tail is at most $2^{-J}$.
This supplies a literal finite expansion with a primitive recursive size.

The elementary logarithm identity
\[
 -\log|1-z|=\mathop{\rm Re}\sum_{j\ge1}z^j/j,
                  \qquad |z|\le1/2,
\]
follows by integrating the finite geometric series for
$z/(1-uz)$ from $u=0$ to $1$. Its tail after $J$ has modulus at most
$2|z|^{J+1}/(J+1)\le2^{-J}$. Apply it to
$z=p^{-\sigma}\exp(-it\log p)$ and add with coefficients $3,4,1$ at
ordinates $0,t,2t$. The summand is
\[
 \frac{p^{-j\sigma}}j
 [3+4\cos(jt\log p)+\cos(2jt\log p)]
 =\frac{2p^{-j\sigma}}j[1+\cos(jt\log p)]^2\ge0.
\]
For the product over $p\le N$, truncation at
$J\ge q+\lceil\log_2N\rceil+10$ has total logarithmic error below
$2^{-q-7}$. The identity $\exp(u+v)=\exp(u)\exp(v)$, proved from its
absolutely convergent power series, now yields the positive product
inequality used below.

One does not have to infer apartness from an exact limiting inequality.
For a direct finite version sufficient for the constants below, suppose
$0<\eta\le1$, choose $B$ with $\eta\ge1/B$, put $q=20+8B$, and use (E5).
Then each of $\zeta(\sigma)$, $\zeta(\sigma+it)$ and
$\zeta(\sigma+2it)$ differs from its finite Euler product by at most
$\varepsilon=2^{-20}B^{-8}$. All six quantities have modulus at most
$4/\eta$. The change in a product of the eight factors, counted with
multiplicity, is at most
\[
 8\varepsilon(4/\eta)^7
 \le2^{17}2^{-20}B^{-8}\eta^{-7}\le1/8.
\]
The finite logarithmic calculation gives its Euler-product side at least
$7/8$ after allowing another explicit $1/8$ approximation error. Hence
the product of the zeta moduli is at least $3/4$, and in particular at
least $1/2$. This weaker numerical version already proves the rounded
$2^{-4}$ lower bound used below. Thus the lower bound has a finite
positive certificate; the argument does not use a principle converting
double-negated nonvanishing into apartness.

\paragraph{The logarithmic derivative series.}
Differentiate the finite Euler product and expand its geometric factors:
\[
 -P_N'(s)/P_N(s)
 =\sum_{p\le N}\frac{(\log p)p^{-s}}{1-p^{-s}}
 =\sum_{p\le N}\sum_{j\ge1}(\log p)p^{-js}.
\]
The unique factorization identification of prime powers gives
\[
 F(s):=-\zeta'(s)/\zeta(s)
              =\sum_{n\ge1}\Lambda(n)n^{-s},\quad\sigma>1. \tag{E6}
\]
This differentiation and limiting step is quantitative. The difference
$\zeta-P_N$ has the omitted-integer expansion used to prove (E5), and
the derivative of that expansion is controlled by the majorant
$\sum_{n>N}(\log n)n^{-\sigma}$; this does not differentiate an
inequality. The reciprocal of $\zeta$ on
$\sigma>1$ has the finite Euler-product bound
$|1/\zeta(s)|\le\prod_p(1+p^{-\sigma})\le\zeta(\sigma)$.
For $c=1+1/b$, $b\ge1$, comparison on each interval $[n-1,n]$ yields
\[
 \sum_{n\ge2}(\log n)n^{-c}
 \le\int_1^\infty\log(2u)u^{-c}du
 =b^2+b\log2\le2b^2.                                     \tag{E7}
\]
The tail after $N$ is at most
$N^{-1/b}[b\log(2N)+b^2]$. For any requested $2^{-q}$, the explicit
choice $N=2^{b^2(q+b+10)^2}$ makes this tail smaller than $2^{-q}$.
To check this, bound the logarithm of the polynomial factor by
$2q+6b+23$ using $\log_2r\le r$, and compare it with
$b(q+b+10)^2$. Enlarging $q$ by the binary length of a preceding
multiplier supplies any desired error budget after multiplication.

\paragraph{A finite scalar Perron kernel.}
For positive rational $v$ and $c\in[1,2]$, define the integral over a
finite vertical segment in the usual parametrized sense. Its integrand
$v^s/(s(s+1))$ has scalar poles at $0,-1$, with residues respectively
$1,-v^{-1}$. We claim
\[
 \frac1{2\pi i}\int_{c-i\infty}^{c+i\infty}
            \frac{v^s}{s(s+1)}ds=(1-v^{-1})_+.              \tag{E8}
\]
Here is a complete finite contour verification. If $v\ge1$, close the
segment on the line $\Re s=-M$, $M\ge2$. Its left edge has absolute
contribution at most $1/(2(M-1))$, because $v^{-M}\le1$ and
\[
 \int_{-\infty}^{\infty}
 \frac{dt}{(M-1)^2+t^2}=\frac\pi{M-1}.
\]
The two horizontal edges at height $U$ have combined contribution at
most $(M+2)\max(1,v^2)/U^2$, using $|s(s+1)|\ge U^2$.
The closed rectangle contains exactly the two displayed scalar poles,
so its integral is $1-v^{-1}$. If $0<v<1$, close instead on
$\Re s=M$, $M>2$. There are no poles; the right-edge and horizontal
bounds are no larger, since $v^M\le1$ and $v^\sigma\le1$ there.
At $v=1$, either closure gives zero, so no apartness from $1$ is needed.
The decision $v\ge1$ is an exact rational comparison.

For a requested error $2^{-q}$ and a rational bound $v\le Y$, the explicit
choices
\[
 M=2^{q+6},\qquad U=2^{q+6}(M+3)(Y+1)^2
\]
make all closing-edge errors below $2^{-q-3}$. The omitted tails on the
original line are bounded by $v^c/(\pi U)$ and satisfy the same bound.
Every contour here is finite and rational. The residue calculation can
be expressed by subtracting the two rational poles and applying the
finite rectangle lemma; the integral of $1/(s-a)$ is evaluated directly
on the four sides by logarithms and arctangents. Alternatively, the same
calculation follows from integrating its elementary derivative away
from $a$ on four rational subrectangles. In either case only the two
explicit real poles occur, and their winding is fixed by the inequalities
$-M<-1<0<c$.

Multiplying (E8) by $y\Lambda(n)$ with $v=y/n$, summing, and using
(E6)--(E7) proves the smoothed Perron formula below. To make every step
finite at accuracy $2^{-q}$, first truncate the Dirichlet series using
the cutoff after (E7), with $q$ enlarged to absorb $3y^2$; then use
(E8) at accuracy $2^{-q}/(1+yN\log N)$ for each of the at most $N$
terms. The order matters: first choose $N$, then the kernel accuracy,
then its auxiliary rectangle sizes. There is no circular choice of
cutoffs. Absolute values and the majorants above account for both
discarded tails. The later PNT contour height $T$ is an independent,
much smaller explicit number.

\section{An explicitly apart zero-free strip}
Put $H=|t|+3$, $L=\log H$, and $a=2^{-88}$. For
$|t|\ge 1/64$, $1-1/(2L)\le\sigma\le2$, Euler summation gives
\[
 \zeta(s)=\sum_{n\le N}n^{-s}+\frac{N^{1-s}}{s-1}
 -s\int_N^\infty\{u\}u^{-s-1}\,du,\qquad N=\lceil H^2\rceil.
\]
Here $L>1$, $\sigma>1/2$, $\log N\le3L$, $N^{1-\sigma}\le4$,
and $|s|\le H$. Consequently
\[
 |\zeta(s)|\le 2^9L,\qquad |\zeta'(s)|\le2^{15}L^2.                 \tag{A1}
\]
For completeness, the finite sum in the first estimate is at most $13L$;
the pole term is at most $256$; the integral term is at most $8/H<3$.
For the derivative, the three contributions are respectively at most
$39L^2$, $768L+16384$, and $8/N+40L/H$. Their sum is below $2^{15}L^2$.
The formula and its derivatives have effective tails: after truncating the
integral at $V$, the zeroth tail is at most $|s|V^{-\sigma}/\sigma$;
the first derivative tail is bounded by
$V^{-\sigma}/\sigma+|s|V^{-\sigma}(\log V/\sigma+1/\sigma^2)$.
The analogous $j$th derivatives have the explicit bound obtained by
integrating $u^{-\sigma-1}(\log u)^j$ by parts $j$ times.
For the bound $N^{1-\sigma}\le4$, use $N\le(10/9)H^2$, hence
$\log N\le(5/2)L$, and $e^{5/4}<4$ (since $e<3$ and $3^5<4^4$).

For $1<\sigma\le2$, the Euler product and
$3+4\cos v+\cos(2v)=2(1+\cos v)^2\ge0$ give
\[
 \zeta(\sigma)^3|\zeta(\sigma+it)|^4
 |\zeta(\sigma+2it)|\ge1.                                      \tag{A2}
\]
This identity requires no choice of a complex logarithm: apply the real
logarithm to the moduli of finite Euler factors, expand their absolutely
convergent real series, and use the displayed nonnegative summands. Pass to
finite approximations with the explicit $\sum_{n>N}n^{-\sigma}$ tail.
Using $\zeta(\sigma)\le2/(\sigma-1)$ and
$|\zeta(\sigma+2it)|\le 2^{10}L$ gives
\[
 |\zeta(\sigma+it)|\ge2^{-4}(\sigma-1)^{3/4}L^{-1/4}.
\]
Set $\delta=aL^{-9}$ and $\sigma_0=1+\delta$.
At $\sigma_0$ this lower bound is $2^{-70}L^{-7}$.
Moving horizontally from $\sigma_0$ by at most $2\delta$ changes $\zeta$
by at most $2^{16}aL^{-7}=2^{-72}L^{-7}$, by (A1).
For $\sigma\ge\sigma_0$ the direct product bound is at least as strong.
Thus, on the whole region
\[
 1-2^{-88}L^{-9}\le\sigma\le2,\quad |t|\ge1/64,
\]
we have the positive, numerical apartness estimates
\[
 |\zeta(s)|\ge2^{-71}L^{-7},\qquad
 \left|\frac{\zeta'(s)}{\zeta(s)}\right|\le2^{86}L^9.             \tag{A3}
\]
In particular this is stronger constructively than merely asserting that
$\zeta(s)\ne0$.

For the finite arithmetization, use the following rational-parameter
version of exactly the same argument. Let $Q$ be a positive integer with
$Q\ge\log(|t|+3)$; take $\delta=2^{-88}Q^{-9}$ and the rational base
point $\sigma_0=1+\delta$. Estimates (A1) hold with $Q$ in place of $L$
on $\sigma\ge1-1/(2Q)$. At $\sigma_0$ the product inequality, even in
its finite $1/2$ version above, gives
$|\zeta(\sigma_0+it)|\ge2^{-4}\delta^{3/4}Q^{-1/4}
=2^{-70}Q^{-7}$. The travel loss is at most $2^{-72}Q^{-7}$.
Thus, throughout $1-2^{-88}Q^{-9}\le\sigma\le2$,
\[
 |\zeta(\sigma+it)|\ge2^{-71}Q^{-7},\qquad
 |\zeta'/\zeta(\sigma+it)|\le2^{86}Q^9.                    \tag{A3r}
\]
The fractional powers in this explanation can be eliminated by fourth
powers: the product lower bound is at least
$2^{-14}\delta^3/Q$, whereas $(2^{-4}\delta^{3/4}Q^{-1/4})^4
=2^{-16}\delta^3/Q$. At the chosen $\delta$, the latter is the fourth
power of the rational number $2^{-70}Q^{-7}$. No irrational sample
point or root-selection operation enters the numerical proof.

\paragraph{The pole neighbourhood.}
Write $\zeta(s)=1/(s-1)+h(s)$, where
$h(s)=1-s\int_1^\infty\{u\}u^{-s-1}\,du$.
If $|\sigma-1|\le1/64$ and $|t|\le1/64$, then $|h|\le4$, $|h'|\le10$,
and $|s-1|\le1/32$. Consequently
$w(s)=(s-1)\zeta(s)=1+(s-1)h(s)$ satisfies $|w|\ge7/8$ and
$|w'|\le5$. Thus
\[
 -\zeta'(s)/\zeta(s)=1/(s-1)+g(s),\qquad |g(s)|\le10,            \tag{A4}
\]
with explicitly computable holomorphic $g=-w'/w$.
For an integer $k\ge2$, set $T=2^k$, $Q=k+2$, $\delta_T=2^{-88}Q^{-9}$,
$\alpha=1-\delta_T$. On the vertical edge $s=\alpha+it$, $|t|\le T$,
(A3)--(A4) and $\log(T+3)\le Q$ imply $|\zeta'/\zeta|\le2^{90}Q^9$.
On the top and bottom horizontal edges, $|\zeta'/\zeta|\le2^{86}Q^9$.
There is exactly the explicitly isolated pole at $s=1$ in the rectangle
$\alpha\le\sigma\le1+1/64$, $|t|\le T$; no zero enumeration is used.

\section{Smoothed Perron and error estimates}
For rational $y\ge e^{64}$ put $R=\log y$. Let $b$ be the least integer
with $y\le2^b$, and put $c=1+1/b$. Then $R\le b\le3R$,
$c\le1+1/64$, and $y^{c+1}\le3y^2$. Define
\[
 \Psi_1(y)=\sum_{n\le y}(y-n)\Lambda(n),\qquad F=-\zeta'/\zeta.
\]
The one-smoothed Perron identity is
\[
 \Psi_1(y)=\frac1{2\pi i}\int_{c-i\infty}^{c+i\infty}
 F(s)\frac{y^{s+1}}{s(s+1)}\,ds.                              \tag{A5}
\]
To prove it, expand $F(s)=\sum\Lambda(n)n^{-s}$ using the absolutely
convergent Euler product and evaluate
$\frac1{2\pi i}\int_{c-i\infty}^{c+i\infty}v^s/(s(s+1))\,ds$
as $(1-v^{-1})_+$, by closing left or right. The only scalar poles are
$0,-1$. At $v=1$ close right; its integral tends to zero like $1/M$ on
$\Re s=M$. When $v\ne1$ is rational, the sign of $v-1$ is decidable and
its distance from $1$ has an explicit rational lower bound. Therefore the
auxiliary rectangle sizes can all be chosen by primitive recursion.
All interchanges are justified by the numerical majorant
$\sum\log(n)n^{-c}\le2b^2\le18R^2$ and $\int(1+t^2)^{-1}dt<4$.

Truncating (A5) at $T$ costs at most $54y^2R^2/T$.
Shift the finite contour to $\alpha=1-\delta_T$. The residue is $y^2/2$.
Since $\int_{-T}^T|s(s+1)|^{-1}dt\le\pi/\alpha\le2\pi$,
the vertical edge costs at most $2^{91}y^2Q^9e^{-\delta_TR}$.
The two horizontal edges cost at most $2^{88}y^2Q^9/T^2$.
Hence
\[
 \left|\Psi_1(y)-\frac{y^2}2\right|
 \le y^2\left(2^{91}Q^9e^{-\delta_TR}
      +2^{88}Q^9/T^2+54R^2/T\right).                       \tag{A6}
\]
The contour step uses the strictly finite rectangle identity proved in Section~\ref{sec:finite-contour-audit}.
Subtract $y^2/(2(s-1))$ from the integrand. In the pole neighbourhood,
write its apparent divided difference as
$\int_0^1 A'(1+u(s-1))du$, with $A(s)=y^{s+1}/(s(s+1))$.
Away from that neighbourhood use (A3). The resulting analytic function
and its first two derivatives have the explicit integer bound
\[
 M=2^{512}(\lceil y\rceil+1)^8(\lceil T\rceil+4)^{40}.
\]
Indeed $|h^{(j)}|\le2^{10}(T+3)$ for $j\le3$, by
$\int_1^\infty(\log u)^ju^{-\sigma-1}du=j!\sigma^{-j-1}$.
For $|t|\ge1/64$, the zeta derivatives through order $3$ are at most
$2^{28}(T+3)$, its reciprocal is at most $2^{71}(T+3)^7$, and the
logarithmic derivative and its first two derivatives are at most
$2^{400}(T+3)^{32}$. For $|t|\le1/64$, use $|w|\ge7/8$ and the
integral divided-difference formula above; derivatives through order $3$
of $A(s)$ are at most $2^{12}y^6$. These bounds are all below $M$.
All contour corners are rational. A rectangular
mesh of side at most $2^{-q}/(100(M+1)(T+1)^2)$, where $M$ bounds the
second derivative, makes the sum of cell-boundary errors below $2^{-q}$.
Each cell uses the first-order Taylor polynomial, whose contour integral
is exactly zero; the $O(M\,\text{side}^3)$ errors sum to the asserted bound.
This is finite telescoping and does not invoke a general compactness theorem.

\section{A convenient all-integer scale}
Let $D_0=2^{100}$ and suppose $k\ge8192$ and
\[
 D_0 k^{10}\le\log_2x<D_0(k+1)^{10}.
\]
Choose $T=2^k$ and use the same $T$ for $y\in[x/2,2x]$.
Then $Q\le2k$, $\delta_T\ge2^{-97}k^{-9}$, and
\[
 \log y\ge2^{98}k^{10},\qquad \log y\le2^{111}k^{10}.
\]
Thus $\delta_T\log y\ge2k$.
For $k\ge8192$, $\log_2k\le k/256$. Here is a wholly integer proof.
Put $j=\lfloor k/256\rfloor\ge32$. At $j=32$ one has
$256(j+1)\le2^j$; the same inequality follows for $j+1$ from
$256(j+2)\le2\cdot256(j+1)$. Quantifier-free induction therefore gives
$k<256(j+1)\le2^j$, whence $k^{256}\le2^{256j}\le2^k$.
Thus the needed polynomial absorptions can be proved without real
differentiation: their $256$th powers are bounded respectively by
$2^{25600+9k}\le2^{32k}$ and
$2^{58368+20k}\le2^{64k}$. Consequently
$2^{100}k^9\le2^{k/8}$ and $54\cdot2^{222}k^{20}\le2^{k/4}$.
The three terms in (A6) are bounded respectively by
$2^{-15k/8}$, $2^{-15k/8}$, and $2^{-3k/4}$.
Their sum is at most $2^{-k/2}$.
Consequently, uniformly for $y\in[x/2,2x]$,
\[
 \left|\Psi_1(y)-y^2/2\right|\le y^2 2^{-k/2}
 \le4x^2 2^{-k/2}.                                         \tag{A7}
\]
Take $h=x2^{-\lfloor k/4\rfloor}$. Since $\psi$ is nondecreasing,
\[
 \frac{\Psi_1(x)-\Psi_1(x-h)}h\le\psi(x)
 \le\frac{\Psi_1(x+h)-\Psi_1(x)}h.
\]
It follows from (A7) that
$|\psi(x)-x|\le h/2+8x^2 2^{-k/2}/h\le9x2^{-k/4}$.
Since $9\le2^{k/8}$, this yields
\[
                    |\psi(x)-x|\le x2^{-k/8}.               \tag{A8}
\]
The $h/2$ estimate is valid because
$h/2\le x2^{-k/4}$ and $2^{\lfloor k/4\rfloor-k/2}\le2^{-k/4}$.
For integer $x$, $k$ is obtained by bounded search from its binary length:
choose the largest $j$ with $D_0 j^{10}\le\lfloor\log_2x\rfloor$.
Thus all rate functions displayed here are primitive recursive.

\paragraph{Status of the analytic estimate.}
The analytic estimates prove an explicit nonoptimal stretched-exponential
PNT error, equivalent to $O(x\exp(-c(\log x)^{1/10}))$ with a positive
explicit $c$. The classical square-root logarithmic rate needs a stronger
zero-free strip and is not proved by these notes. An iPRA formalization
still has to instantiate the indicated finite rational complex-arithmetic
identities and their error budgets; the fact that the numerical bounds are
primitive recursive is not, by itself, a proof that the argument is in iPRA.
